\documentclass[
    12pt,
    a4paper
]{report}



\usepackage{graphicx}
\usepackage{geometry}

\title{ELEC473 Assignment 4 - Synthesising the NIOS II Processor}
\date{\today}
\author{Ryan McKee - 201954749}


% TODO When report complete

\begin{document}

\maketitle

\begin{abstract}
    
\end{abstract}


\chapter{Introduction}

We focus on expanding familiarity of FPGA \(Field Programmable Gateway\) engineering through 
this assignments focus on \bold{QSYS}~\cite{qsys} and the \bold{NIOS-II}~\cite{nios_ii} processor. 
The NIOS-II being \bold{Altera's}~\cite{altera} own 32-bit embedded processor 
mirco-architecture specifically designed for syntehesis on FPGA, The NIOS-II is a successor 
of Altera's first configurable 16-bit embedded processor Nios, introduced in 2000. During this assignment 
we utilize the Cyclone-II DE2 FPGA which we syntehsis designs, bench test and design specifically for during the course
of our work.


The NIOS-II hosts the following design parameters: 
-   32 general-purpose 32-bit registers
-   Full 32-bit Instruction set, data path, and address space.
- Singla-instruction 32 x 32 mulitply and divide producing a 32-bit result. 

This soft-core \(A soft microprocessor is a microprocessor core that can be wholly implemented using logic syntehsis. It can be implemented 
via different semiconductor devices containing Programmable logic (e.g. FPGA, CPLD) including bothhigh-end and commodity validations~\cite{soft-microprocessor}\) nature of NIOS-II enables hardware engineers to specify and generate 
custom NIOS II core, tailed for their own cusom application requirements. 

Qsys ( now known as platform designer) is a system integration tool for Intel Quartus Prime.
It helps connect multiple intellectual property (IP) blocks, 
processors (like Nios II), and memory interfaces using standard interconnects like Avalon and AXI.
It can be utilized for generated HDL.

This assignment sees us in Part A sythesise a NIOS II processor with 20KB of on-chip memory, a timer, a JTAG UART,
8 parallel I/O pins and system ID component. part a sees us interface the 512 KB SRAM and 8 MB SDRAM on the Altera DE2 Board 
to this design and test it using memory test programs within NIOS II IDE we will show Name and ID numbers
in terminal each time the memory is testing. 

part b continues on part a requiring us to develop a custom instruction to count the number of leading 1s in the 32 bit number passed into 
the instruction. and write a custom program to the this custom instruction. I will also develop a test routine in C or assembler utilizing compiler that
has the same instruction fucntionality implemented to compare the speed of the custom instruction against software implementation. 

Part C we shall utilize PWM to control the inensite of of LEDR10. 


\chapter{Part A - SRAM \& SDRAM}



\chapter{Part B - Custom Instruction}

\chapter{Part C - PWM Module}

\bibliographystyle{plain}
\bibliography{bibliography}

\end{document}